\Needspace{8\baselineskip}\originalchapter{平衡与稳定的数值工作流}\label{ch:numerics}
\originalsection{先确定所解模型}
数值工作先固定场、通量、方向、无量纲参数和边界, 再选离散法. 将不同讲义里不同归一的 \ensuremath{\Psi} 和 beta 混用, 程序即使收敛也会解错方程. 对 GS 可用有限差分/有限元求直接场, 或反坐标矩投影求有限参数磁面. 前者需要网格边界和离散残差, 后者需要矩方程、Jacobians和投影外残差; 两者的``误差''含义不同.

本册编程试卷实际附 Python 代码、CSV和核验 JSON. 源码包中的数值代码只用标准库及明确列出的 NumPy/SciPy/Matplotlib, 固定随机性或完全确定求积. 所绘图从真实输出数据生成, 不是为匹配预期而手工捏造点.

\originalsection{求积和根的独立检查}
面积积分先写 Jacobian, 在磁轴、分离面附近分析奇性, 再用分段角求积或适当变量代换. 参数根求解需检查约束、不同初值、残差以及加密后根是否保持. 坐标嵌套要求 Jacobian 不变号; q 积分分母若趋零, 精度加密也不能把物理发散压成有限数.

对稳定矩阵, 先验证对称性和已知低压圆柱解析极限, 再求特征值. 壁规则若随参数突然把某壁移至无穷, 势能本身有跳变, 不能把跨跳变的数值二分结果叫做零特征值的 marginal point. 求积误差、模型分支跳变和真实过零应分别列出.

\originalsection{报告的量词和误差}
真实计算应报告输入值、无量纲化、方程残差、网格/求积加密、分支条件和输出单位. 在三谐波空间 \ensuremath{\delta}W\ensuremath{\ge}0只称``该空间未发现负方向''; 若高 m 扫描也有限, 说明最大 m 和参数范围. 不将一个模型的上确界叫作装置全物理稳定上限. 这使课程型算例能供自学复现, 同时保留物理判断的边界.
